% Hafta 5 — -keep dengesi
% Derleme: py -3.12 tools/latex/derle.py docs/week-5/assets/h05-09-keep-dengesi.tex
\documentclass[tikz,border=8pt]{standalone}
\usepackage{cen429-cizim}
\begin{document}
\begin{tikzpicture}
  \node[baslik] (b) at (0,0) {-keep dengesi};
  \node[altbaslik] at (b.south west) {doğrusu: yalnız dışarıdan erişilen giriş noktalarını keep et};
  \node[kotukutu=70mm, anchor=north west] (sol) at (0,-2.0)
    {\kb{ÇOK keep}\\[2mm]
     reflection, serileştirme,\\tüm public API korunur\\[3mm]
     $\rightarrow$ adlar değişmez\\$\rightarrow$ gizleme ANLAMSIZ};
  \node[uyarikutu=70mm, right=14mm of sol] (sag)
    {\kb{AZ keep}\\[2mm]
     yalnız dış giriş noktaları\\[3mm]
     $\rightarrow$ adlar kısalır\\$\rightarrow$ ama reflection KIRILIR\\(çalışma anı hatası)};
  \node[fit=(sol)(sag), inner sep=0pt] (grp) {};
  \node[dip, text width=155mm, anchor=north west] at ($(grp.south west)+(0,-8mm)$)
    {Her -keep kuralı bir gizleme boşluğudur; en dar kuralı yazın ve test edin.};
\end{tikzpicture}
\end{document}
